%==========================================
%
% SIBGRAPI 2026 paper template
% Example of IEEEtran.cls
%
%==========================================

% *** Authors should verify (and, if needed, correct) their LaTeX system  ***
% *** with the testflow diagnostic prior to trusting their LaTeX platform ***
% *** with production work. The IEEE's font choices and paper sizes can   ***
% *** trigger bugs that do not appear when using other class files.       ***                          ***
% The testflow support page is at:
% http://www.michaelshell.org/tex/testflow/

\documentclass[10pt,conference]{IEEEtran}


% Some very useful LaTeX packages include:
% (uncomment the ones you want to load)


% *** MISC UTILITY PACKAGES ***
%
%\usepackage{ifpdf}
% Heiko Oberdiek's ifpdf.sty is very useful if you need conditional
% compilation based on whether the output is pdf or dvi.
% usage:
% \ifpdf
%   % pdf code
% \else
%   % dvi code
% \fi
% The latest version of ifpdf.sty can be obtained from:
% http://www.ctan.org/pkg/ifpdf
% Also, note that IEEEtran.cls V1.7 and later provides a builtin
% \ifCLASSINFOpdf conditional that works the same way.
% When switching from latex to pdflatex and vice-versa, the compiler may
% have to be run twice to clear warning/error messages.






% *** CITATION PACKAGES ***
%
\usepackage{cite}
% cite.sty was written by Donald Arseneau
% V1.6 and later of IEEEtran pre-defines the format of the cite.sty package
% \cite{} output to follow that of the IEEE. Loading the cite package will
% result in citation numbers being automatically sorted and properly
% "compressed/ranged". e.g., [1], [9], [2], [7], [5], [6] without using
% cite.sty will become [1], [2], [5]--[7], [9] using cite.sty. cite.sty's
% \cite will automatically add leading space, if needed. Use cite.sty's
% noadjust option (cite.sty V3.8 and later) if you want to turn this off
% such as if a citation ever needs to be enclosed in parenthesis.
% cite.sty is already installed on most LaTeX systems. Be sure and use
% version 5.0 (2009-03-20) and later if using hyperref.sty.
% The latest version can be obtained at:
% http://www.ctan.org/pkg/cite
% The documentation is contained in the cite.sty file itself.






% *** GRAPHICS RELATED PACKAGES ***
%
\ifCLASSINFOpdf
   \usepackage[pdftex]{graphicx}
  % declare the path(s) where your graphic files are
   \graphicspath{{figs/}}
  % and their extensions so you won't have to specify these with
  % every instance of \includegraphics
   \DeclareGraphicsExtensions{.pdf,.jpeg,.png}
\else
  % or other class option (dvipsone, dvipdf, if not using dvips). graphicx
  % will default to the driver specified in the system graphics.cfg if no
  % driver is specified.
   \usepackage[dvips]{graphicx}
  % declare the path(s) where your graphic files are
   \graphicspath{{../figs/}}
  % and their extensions so you won't have to specify these with
  % every instance of \includegraphics
   \DeclareGraphicsExtensions{.eps}
\fi
% graphicx was written by David Carlisle and Sebastian Rahtz. It is
% required if you want graphics, photos, etc. graphicx.sty is already
% installed on most LaTeX systems. The latest version and documentation
% can be obtained at: 
% http://www.ctan.org/pkg/graphicx
% Another good source of documentation is "Using Imported Graphics in
% LaTeX2e" by Keith Reckdahl which can be found at:
% http://www.ctan.org/pkg/epslatex
%
% latex, and pdflatex in dvi mode, support graphics in encapsulated
% postscript (.eps) format. pdflatex in pdf mode supports graphics
% in .pdf, .jpeg, .png and .mps (metapost) formats. Users should ensure
% that all non-photo figures use a vector format (.eps, .pdf, .mps) and
% not a bitmapped formats (.jpeg, .png). The IEEE frowns on bitmapped formats
% which can result in "jaggedy"/blurry rendering of lines and letters as
% well as large increases in file sizes.
%
% You can find documentation about the pdfTeX application at:
% http://www.tug.org/applications/pdftex





% *** MATH PACKAGES ***
%
\usepackage[cmex10]{amsmath}
% A popular package from the American Mathematical Society that provides
% many useful and powerful commands for dealing with mathematics.
%
% Note that the amsmath package sets \interdisplaylinepenalty to 10000
% thus preventing page breaks from occurring within multiline equations. Use:
\interdisplaylinepenalty=2500
% after loading amsmath to restore such page breaks as IEEEtran.cls normally
% does. amsmath.sty is already installed on most LaTeX systems. The latest
% version and documentation can be obtained at:
% http://www.ctan.org/pkg/amsmath
\usepackage{amsthm}
\newtheorem{definition}{Definition}




% *** SPECIALIZED LIST PACKAGES ***
%
\usepackage{algorithmic}
% algorithmic.sty was written by Peter Williams and Rogerio Brito.
% This package provides an algorithmic environment fo describing algorithms.
% You can use the algorithmic environment in-text or within a figure
% environment to provide for a floating algorithm. Do NOT use the algorithm
% floating environment provided by algorithm.sty (by the same authors) or
% algorithm2e.sty (by Christophe Fiorio) as the IEEE does not use dedicated
% algorithm float types and packages that provide these will not provide
% correct IEEE style captions. The latest version and documentation of
% algorithmic.sty can be obtained at:
% http://www.ctan.org/pkg/algorithms
% Also of interest may be the (relatively newer and more customizable)
% algorithmicx.sty package by Szasz Janos:
% http://www.ctan.org/pkg/algorithmicx




% *** ALIGNMENT PACKAGES ***
%
\usepackage{array}
\usepackage{tikz}
\usepackage{svg}
\usetikzlibrary{arrows.meta,positioning,shapes.geometric}
% Frank Mittelbach's and David Carlisle's array.sty patches and improves
% the standard LaTeX2e array and tabular environments to provide better
% appearance and additional user controls. As the default LaTeX2e table
% generation code is lacking to the point of almost being broken with
% respect to the quality of the end results, all users are strongly
% advised to use an enhanced (at the very least that provided by array.sty)
% set of table tools. array.sty is already installed on most systems. The
% latest version and documentation can be obtained at:
% http://www.ctan.org/pkg/array


% IEEEtran contains the IEEEeqnarray family of commands that can be used to
% generate multiline equations as well as matrices, tables, etc., of high
% quality.




% *** SUBFIGURE PACKAGES ***
\ifCLASSOPTIONcompsoc
  \usepackage[caption=false,font=normalsize,labelfont=sf,textfont=sf]{subfig}
\else
  \usepackage[caption=false,font=footnotesize]{subfig}
\fi
% subfig.sty, written by Steven Douglas Cochran, is the modern replacement
% for subfigure.sty, the latter of which is no longer maintained and is
% incompatible with some LaTeX packages including fixltx2e. However,
% subfig.sty requires and automatically loads Axel Sommerfeldt's caption.sty
% which will override IEEEtran.cls' handling of captions and this will result
% in non-IEEE style figure/table captions. To prevent this problem, be sure
% and invoke subfig.sty's "caption=false" package option (available since
% subfig.sty version 1.3, 2005/06/28) as this is will preserve IEEEtran.cls
% handling of captions.
% Note that the Computer Society format requires a larger sans serif font
% than the serif footnote size font used in traditional IEEE formatting
% and thus the need to invoke different subfig.sty package options depending
% on whether compsoc mode has been enabled.
%
% The latest version and documentation of subfig.sty can be obtained at:
% http://www.ctan.org/pkg/subfig




% *** FLOAT PACKAGES ***
%
%\usepackage{fixltx2e}
% fixltx2e, the successor to the earlier fix2col.sty, was written by
% Frank Mittelbach and David Carlisle. This package corrects a few problems
% in the LaTeX2e kernel, the most notable of which is that in current
% LaTeX2e releases, the ordering of single and double column floats is not
% guaranteed to be preserved. Thus, an unpatched LaTeX2e can allow a
% single column figure to be placed prior to an earlier double column
% figure.
% Be aware that LaTeX2e kernels dated 2015 and later have fixltx2e.sty's
% corrections already built into the system in which case a warning will
% be issued if an attempt is made to load fixltx2e.sty as it is no longer
% needed.
% The latest version and documentation can be found at:
% http://www.ctan.org/pkg/fixltx2e


%\usepackage{stfloats}
% stfloats.sty was written by Sigitas Tolusis. This package gives LaTeX2e
% the ability to do double column floats at the bottom of the page as well
% as the top. (e.g., "\begin{figure*}[!b]" is not normally possible in
% LaTeX2e). It also provides a command:
%\fnbelowfloat
% to enable the placement of footnotes below bottom floats (the standard
% LaTeX2e kernel puts them above bottom floats). This is an invasive package
% which rewrites many portions of the LaTeX2e float routines. It may not work
% with other packages that modify the LaTeX2e float routines. The latest
% version and documentation can be obtained at:
% http://www.ctan.org/pkg/stfloats
% Do not use the stfloats baselinefloat ability as the IEEE does not allow
% \baselineskip to stretch. Authors submitting work to the IEEE should note
% that the IEEE rarely uses double column equations and that authors should try
% to avoid such use. Do not be tempted to use the cuted.sty or midfloat.sty
% packages (also by Sigitas Tolusis) as the IEEE does not format its papers in
% such ways.
% Do not attempt to use stfloats with fixltx2e as they are incompatible.
% Instead, use Morten Hogholm'a dblfloatfix which combines the features
% of both fixltx2e and stfloats:
%
% \usepackage{dblfloatfix}
% The latest version can be found at:
% http://www.ctan.org/pkg/dblfloatfix




% *** PDF, URL AND HYPERLINK PACKAGES ***
%
\usepackage{url}
% url.sty was written by Donald Arseneau. It provides better support for
% handling and breaking URLs. url.sty is already installed on most LaTeX
% systems. The latest version and documentation can be obtained at:
% http://www.ctan.org/pkg/url
% Basically, \url{my_url_here}.




% *** Do not adjust lengths that control margins, column widths, etc. ***
% *** Do not use packages that alter fonts (such as pslatex).         ***
% There should be no need to do such things with IEEEtran.cls V1.6 and later.
% (Unless specifically asked to do so by the journal or conference you plan
% to submit to, of course. )


% correct bad hyphenation here
\hyphenation{op-tical net-works semi-conduc-tor}

%------------------------------------------------------------------------- 
% change the % on next lines to produce the final camera-ready version 
\newif\iffinal
\finalfalse
%\finaltrue
\newcommand{\cmtid}{115}
%------------------------------------------------------------------------- 

\iffinal
\else
\usepackage[switch]{lineno}
\fi

\begin{document}
%
% paper title
% Titles are generally capitalized except for words such as a, an, and, as,
% at, but, by, for, in, nor, of, on, or, the, to and up, which are usually
% not capitalized unless they are the first or last word of the title.
% Linebreaks \\ can be used within to get better formatting as desired.
% Do not put math or special symbols in the title.
\title{Resolution-Bounded KDE: Real-Time Density Heatmaps for Large Weighted Timeseries}


% author names and affiliations
% use a multiple column layout for up to two different
% affiliations

\iffinal

% author names and affiliations
% use a multiple column layout for up to three different
% affiliations
\author{\IEEEauthorblockN{Gabriel Eduardo de Lima Machado}
\IEEEauthorblockA{Institute of Informatics \\
Federal University of Rio Grande do Sul\\
Porto Alegre, RS, Brazil\\
Email: gabriel.delmachado@gmail.com}
\and
\IEEEauthorblockN{Carla M. Dal Sasso Freitas}
\IEEEauthorblockA{Institute of Informatics \\
Federal University of Rio Grande do Sul\\
Porto Alegre, RS, Brazil\\
Email: carla@inf.ufrgs.br}
}

% conference papers do not typically use \thanks and this command
% is locked out in conference mode. If really needed, such as for
% the acknowledgment of grants, issue a \IEEEoverridecommandlockouts
% after \documentclass

% for over three affiliations, or if they all won't fit within the width
% of the page, use this alternative format:
% 
%\author{\IEEEauthorblockN{Michael Shell\IEEEauthorrefmark{1},
%Homer Simpson\IEEEauthorrefmark{2},
%James Kirk\IEEEauthorrefmark{3}, 
%Montgomery Scott\IEEEauthorrefmark{3} and
%Eldon Tyrell\IEEEauthorrefmark{4}}
%\IEEEauthorblockA{\IEEEauthorrefmark{1}School of Electrical and Computer Engineering\\
%Georgia Institute of Technology,
%Atlanta, Georgia 30332--0250\\ Email: see http://www.michaelshell.org/contact.html}
%\IEEEauthorblockA{\IEEEauthorrefmark{2}Twentieth Century Fox, Springfield, USA\\
%Email: homer@thesimpsons.com}
%\IEEEauthorblockA{\IEEEauthorrefmark{3}Starfleet Academy, San Francisco, California 96678-2391\\
%Telephone: (800) 555--1212, Fax: (888) 555--1212}
%\IEEEauthorblockA{\IEEEauthorrefmark{4}Tyrell Inc., 123 Replicant Street, Los Angeles, California 90210--4321}}

\else
  \author{SIBGRAPI Paper ID: \cmtid \\ }
  \linenumbers
\fi


% make the title area
\maketitle

% As a general rule, do not put math, special symbols or citations
% in the abstract
\begin{abstract}
A timeseries is usually shown as a time-versus-value plot, leaving any third scalar (for example, traded volume in a temporal series of stock trading) underused. In our work, we encode this third variable as weighted Kernel Density Estimation (KDE) and render it as a heatmap, so all three dimensions are visible in one chart. At hundreds of thousands of points, CPU KDE and na\"{i}ve per-point GPU splatting miss real-time pan/zoom budgets. We therefore merge nearby points in screen space into weighted kernels before rendering, reducing the number of evaluations to roughly the screen resolution. The resulting three-stage pipeline --- \textit{accumulation}, \textit{reduction}, and \textit{dynamic color mapping} --- uses instanced drawing for Gaussian splats and produces perceptually calibrated density maps at interactive rates without server-side pre-aggregation.
\end{abstract}

% no keywords




% For peerreview papers, this IEEEtran command inserts a page break and
% creates the second title. It will be ignored for other modes.
\IEEEpeerreviewmaketitle



%==============================
% Instruções de submissão em https://sibgrapi.sbc.org.br/2026/#/submission
%==============================

\section{Introduction}
\label{sec:motivation}

Visual representations, mostly charts, have been used for centuries to depict temporal data, helping researchers understand and illustrate facts and phenomena. Several examples of time-series visualizations, including the classic ones by William Playfair, can be easily found in information visualization books, especially in Aigner et al.'s \cite{Aigner23TimeViz2nd}. 

%A timeseries is naturally plotted as time versus value, but many datasets also include a third scalar (for example traded volume). We encode this third value as weighted KDE: each sample contributes a weighted Gaussian in chart space, and color encodes accumulated density. Prior work most closely related to this formulation is summarized in Section~\ref{sec:related-work}.

Temporal data may have multiple values per time point, and over the years, several visualization techniques have been developed, beyond the popular time-versus-value plot. However, when the problem at hand requires observing a large time series in real time, such complex visualizations do not scale. Even adding a second value to a time point requires a different approach. 

We propose encoding this second value as weighted Kernel Density Estimation: each sample contributes a weighted Gaussian in chart space, and color encodes accumulated density. %Prior work most closely related to this formulation is summarized in Section~\ref{sec:related-work}. 

The challenge is scale. With hundreds of thousands of points, CPU evaluation costs $O(N \times W \times H)$ and is too slow for interactive pan/zoom. Na\"{i}ve GPU splatting (one kernel quad per point) reduces constants but still suffers heavy overdraw \cite{10.1109/TVCG.2013.65}; with strong overlap, fragment work can reach billions of evaluations per frame.

We address this with a screen-space kernel-merging approximation: nearby points are replaced by one weighted representative, and merging is recomputed after each view change. This bounds submitted kernels by display resolution rather than dataset size while keeping visual error low when the threshold is near the kernel radius. The resulting accumulation/reduction/dynamic-color-mapping pipeline supports interactive, perceptually calibrated density maps (see Figure \ref{fig:bitcoin-sample}).

\section{Related Work}
\label{sec:related-work}

Our approach is most closely related to the interactive KDE \cite{5742387} and Gaussian-mixture rendering \cite{article} lines of work, combining the interactive KDE perspective for exploratory analysis with Gaussian-mixture rendering ideas from large-scale particle visualization, while targeting weighted financial timeseries heatmaps and screen-space merge control. We briefly review works from these two lines of research in this section.

Lampe and Hauser \cite{5742387} introduced an interactive Kernel Density Estimation (KDE) framework for streaming data visualization, including viewport-adaptive bandwidth selection and a GPU implementation for interactive updates under pan and zoom. Their work established a practical foundation for coupling KDE with interactive visual analytics.

More recently, Peng et al. \cite{article} presented a Gaussian mixture model-based splatting pipeline for large-scale particle simulations, where Gaussian components are rendered directly from a compressed representation to avoid expensive reconstruction. That work demonstrates the viability of Gaussian-mixture primitives for high-throughput visualization and informs our use of Gaussian-based accumulation in a performance-critical pipeline.

A second line of work addresses the overdraw that arises when many points are plotted in the same screen region. Bachthaler and Weiskopf's continuous scatterplots \cite{4658159} recast discrete point plots as a continuous density field, while Splatterplots \cite{10.1109/TVCG.2013.65} bound the density of marks drawn per unit of screen space and leverage the GPU to scale to millions of points. These methods share our central premise: the displayed resolution, rather than the raw record count, should govern how much is drawn. We adopt the same principle but target a weighted KDE density field rather than abstracted contours or outlier glyphs.

The notion that interactive scalability should be limited by display resolution is also central to scalable aggregation systems. imMens \cite{10.1111/cgf.12129} precomputes binned aggregates and renders them on the GPU for frame-rate brushing and linking over billions of records, and Nanocubes \cite{10.1109/TVCG.2013.179} maintain an in-memory data cube that answers aggregate queries with bounded screen error across scales. These systems rely on precomputed, server- or memory-resident aggregations; in contrast, our screen-space merge is recomputed per view directly on the visible subset, avoiding any pre-aggregation step.

Finally, our accumulation stage is methodologically a splatting technique: each kernel is rasterized as an anisotropic Gaussian footprint defined by a covariance matrix. This connects our pipeline to recent 3D Gaussian splatting \cite{10.1145/3592433}, which renders scenes as large collections of anisotropic Gaussians with optimized covariances at real-time rates. We borrow the same instanced, covariance-driven splatting formulation, but apply it to 2D weighted density estimation rather than radiance-field synthesis.

%Our approach is most closely related to the interactive KDE \cite{5742387} and Gaussian-mixture rendering \cite{article} lines of work, combining the interactive KDE perspective for exploratory analysis with Gaussian-mixture rendering ideas from large-scale particle visualization, while targeting weighted financial timeseries heatmaps and screen-space merge control.


\section{Overview}
\label{sec:data}

\subsection{Data Preparation}

Our target application is the visual exploration of large series of financial data, mainly trading series. Each data point carries three values: a timestamp, a price level, and a scalar weight. For financial tick data, weight is naturally derived from traded volume, making denser activity periods visually brighter regardless of the number of individual events.

%Before rendering, timestamps are normalized to a unit interval. This ensures that coordinate arithmetic in subsequent stages operates on well-conditioned values regardless of the absolute time range, preventing precision loss at extreme zoom levels. The original scale factors are retained separately for axis label reconstruction.

All examples in this paper use the same synthetic dataset generation procedure: a random-walk trajectory for the signal, combined with the same amplitude function to define the distribution of quantities (Figure~\ref{fig:synthetic-amplitude}), with Gaussian noise added. The only exception is the Bitcoin chart (Fig.~\ref{fig:bitcoin-sample}), which is built from real market data.

Figure~\ref{fig:bitcoin-sample} shows a representative dataset slice used in our experiments, where long-horizon price motion and localized bursts of trading activity coexist and motivate density-based rendering.

\begin{figure}[t]
  \centering
  \includegraphics[width=1\linewidth]{figs/bitcoin-negotiation.png}
  \caption{Bitcoin trading sample from Coinbase, with 100 thousand samples of 5 minutes compressed into 205 Gaussian kernels, highlighting the February downturn and unusual volume concentration around Feb 05 where more than 390 thousand bitcoins were traded.}
  \label{fig:bitcoin-sample}
\end{figure}


\begin{figure}[t]
  \centering
  \includesvg[width=1\linewidth]{figs/quantity_amplitude}
  \caption{Amplitude function used to define the quantity distribution in the synthetic data generation process shared by all non-Bitcoin examples.}
  \label{fig:synthetic-amplitude}
\end{figure}


Before rendering, timestamps are normalized to a unit interval. This ensures that coordinate arithmetic in subsequent stages operates on well-conditioned values regardless of the absolute time range, preventing precision loss at extreme zoom levels. The original scale factors are retained separately for axis label reconstruction.


%\section{Adaptive Downsampling}
\subsection{Adaptive Downsampling}
\label{sec:downsampling}

%\subsection{View-Range Clipping}

Because data is sorted by time, only points falling within the current view window need to be considered. Binary search extracts this subset in logarithmic time, keeping subsequent work proportional to visible data rather than total data.

\subsubsection{Weighted Merge}

Visible points are downsampled by a \textit{weighted merge} pass before GPU upload. Scanning left-to-right, the algorithm maintains one open cluster and tests whether the next point can be absorbed while keeping both the earliest point and the candidate within a screen-space threshold of the updated weighted centroid. If not, it flushes the cluster and starts a new one.

Each merged point keeps the sum of constituent weights, preserving the total signal. The earliest-point check prevents centroid drift from slowly violating the threshold. Since points are time-sorted, each point is visited once with no backtracking, yielding $O(N)$ time over the visible subset. The merge is recomputed on each pan/zoom event to keep transitions smooth.

Figure~\ref{fig:merge-downsample} illustrates the exact decision used in our implementation. The left side shows the geometric test performed for each incoming point: after computing the hypothetical weighted centroid $\mu'$, the algorithm accepts the candidate only if both the first point of the open cluster and the new point remain within the threshold radius $\tau$ in pixel space. The right side shows the resulting one-pass control flow.

\begin{figure}[t]
  \centering
  \begin{minipage}[t]{\linewidth}
    \vspace{0pt}
    \centering
    \begin{tikzpicture}[
      font=\scriptsize,
      >=Latex,
      line width=0.45pt,
      flow/.style={draw, rounded corners=1.5pt, align=center, minimum width=18mm, minimum height=6mm, text width=20mm},
      decision/.style={draw, diamond, aspect=1.65, align=center, inner xsep=0.8mm, inner ysep=0.6mm, text width=16mm}
    ]
      \node[anchor=west, font=\bfseries\footnotesize] at (-0.1,3.1) {One-pass merge loop};

      \node[flow] (start) at (2.4,2.45) {Read next point $p_i$};
      \node[flow] (centroid) at (2.4,1.05) {Compute hypothetical\\centroid $\mu'$};
      \node[decision] (inside) at (2.4,-0.55) {All inside\\$\tau$?};
      \node[flow] (flush) at (0.25,-2.2) {Flush previous cluster\\Start new one with $p_i$};
      \node[flow] (absorb) at (4.55,-2.2) {Absorb point\\update centroid, weight};

      \draw[->] (start) -- (centroid);
      \draw[->] (centroid) -- (inside);
      \draw[->] (inside) -- node[above] {no} (flush);
      \draw[->] (inside) -- node[above] {yes} (absorb);
      \draw[->] (absorb.north) |- (start.east);
      \draw[->] (flush.north) |- (start.west);
    \end{tikzpicture}
  \end{minipage}
  \caption{Diagram of the weighted-merge downsampling pass implemented by \texttt{mergeDownsample}. Left: geometric acceptance criterion in screen space. Right: control flow of the single left-to-right scan. The open cluster is updated only when the first point and the incoming point both remain within threshold $\tau$ of the hypothetical new centroid.}
  \label{fig:merge-downsample}
\end{figure}

\subsubsection{Why Shape-Preserving Simplification Is Insufficient}

Classical timeseries simplification algorithms such as Largest-Triangle-Three-Buckets (LTTB) \cite{Steinarsson2013} and Ramer-Douglas-Peucker (RDP) \cite{Douglas1973,10589058} are designed to preserve the visual shape of a line chart: they select or retain points whose positions best reconstruct the geometry of the original polyline, discarding points they deem geometrically redundant. This objective is appropriate for a line chart, where the eye follows the trajectory of the series. It is, however, the wrong objective for KDE.

In a density heatmap the quantity that must be conserved is the total weight --- the accumulated momentum of the signal --- not its geometric shape. A point that is geometrically redundant (lying close to a chord between its neighbors) may nonetheless carry substantial weight that, if discarded, would leave the density field darker than the data warrants. Conversely, a point selected by LTTB because it lies far from the chord may carry negligible weight, contributing little to the density while consuming a kernel evaluation. Both cases produce a density estimate that misrepresents the underlying distribution.

The weighted merge strategy addresses this directly: the sole criterion for collapsing points is screen-space proximity, and the total weight is always conserved in the merged centroid. No information about the signal's magnitude is lost; only the spatial resolution, which is already below the pixel threshold and therefore invisible.


\section{Mathematical Basis for Gaussian Merging}
\label{sec:gaussian-merging}

This section formalizes our merge operation and relates it to the Gaussian-mixture reduction (GMR) literature from target tracking, where a mixture is approximated by one with fewer components by successively merging pairs \cite{Salmond2009,Williams2003,4383588}; Crouse et al. \cite{inproceedings} survey and compare these algorithms. Our moment-preserving merge is the same second-order moment match used as the per-pair reduction step in that work, specialized here to screen-space proximity rather than the KL-based pair-selection criterion of Runnalls \cite{4383588}.

\begin{figure}[t]
  \centering
  \includegraphics[width=1\linewidth]{figs/gaussian_merge2.png}
  \caption{Two-component Gaussian merge example for equal-variance kernels. At centroid separation $\rho=\lVert\overrightarrow{\boldsymbol{\Delta}}\rVert/\sigma=1$, the moment-preserving single-Gaussian approximation induces low integrated squared error.}
  \label{fig:two-gaussian-merge}
\end{figure}

At the component level, we use the full matrix-form weighted Gaussian
\begin{equation}
\mathcal{N}(\overrightarrow{\boldsymbol{x}}; A, \overrightarrow{\boldsymbol{\mu}}, \overrightarrow{\mathbf{P}}) = A\,\exp\!\left(-\frac{1}{2}(\overrightarrow{\boldsymbol{x}}-\overrightarrow{\boldsymbol{\mu}})^T \overrightarrow{\mathbf{P}}^{-1}(\overrightarrow{\boldsymbol{x}}-\overrightarrow{\boldsymbol{\mu}})\right),
\end{equation}
with amplitude defined from the total weight $w$ as
\begin{equation}
A(w)=\frac{w}{\sqrt{(2\pi)^d\det(\overrightarrow{\mathbf{P}})}},
\end{equation}
where $\overrightarrow{\boldsymbol{x}}\in\mathbb{R}^d$, $\overrightarrow{\boldsymbol{\mu}}\in\mathbb{R}^d$, and $\overrightarrow{\mathbf{P}}\in\mathbb{R}^{d\times d}$ is symmetric positive definite.

Each visible sample is modeled as one Gaussian-mixture component
\begin{equation}
\left(w_i, \overrightarrow{\boldsymbol{\mu}}_i, \overrightarrow{\mathbf{P}}_i\right),
\end{equation}
where $w_i$ is its scalar weight (volume), $\overrightarrow{\boldsymbol{\mu}}_i \in \mathbb{R}^2$ is its chart-space position, and $\overrightarrow{\mathbf{P}}_i$ is its kernel covariance.

For a set of $N$ components, the moment-preserving single-Gaussian approximation is
\begin{equation}
w = \sum_{i=1}^{N} w_i,
\end{equation}
\begin{equation}
\overrightarrow{\boldsymbol{\mu}} = \frac{1}{w}\sum_{i=1}^{N} w_i\overrightarrow{\boldsymbol{\mu}}_i,
\end{equation}
\begin{equation}
\overrightarrow{\mathbf{P}} = \frac{1}{w}\sum_{i=1}^{N} w_i\left[\overrightarrow{\mathbf{P}}_i + (\overrightarrow{\boldsymbol{\mu}}_i-\overrightarrow{\boldsymbol{\mu}})(\overrightarrow{\boldsymbol{\mu}}_i-\overrightarrow{\boldsymbol{\mu}})^T\right].
\end{equation}

For the merged Gaussian component, the peak amplitude in 2D can be written from its total weight and covariance determinant as
\begin{equation}
A = \frac{w}{2\pi\sqrt{\det(\overrightarrow{\mathbf{P}})}},
\end{equation}
which follows from the normalization constant of a weighted Gaussian density.

\section{GPU Rendering Pipeline}
\label{sec:pipeline}

The rendering pipeline consists of three GPU passes submitted together per frame. In Pass 1, Gaussian kernels are rasterized with instanced drawing (one instance per merged kernel) from a single quad mesh.

\begin{figure}[t]
  \centering
  \begin{tikzpicture}[
    font=\footnotesize,
    node distance=4.5mm,
    stage/.style={draw, rounded corners=1.5pt, align=center, minimum height=9mm, text width=0.56\linewidth},
    io/.style={draw, align=center, minimum height=8mm, text width=0.56\linewidth},
    arrow/.style={-Latex, line width=0.5pt}
  ]
    \node[io] (raw) {Raw points $(t,p,w)$};
    \node[stage, below=of raw] (merge) {CPU weighted merge ($\tau$ in pixel space)};
    \node[io, below=of merge] (inst) {Merged kernels (instance buffer upload)};
    \node[stage, below=of inst] (vs) {Pass 1a: Vertex shader (instanced screen-aligned quads, one instance per kernel)};
    \node[stage, below=of vs] (acc) {Pass 1b: Fragment shader accumulation (quad rasterization + additive blending into HDR KDE texture)};
    \node[stage, below=of acc] (red) {Pass 2: Compute shader reduction (parallel max to compute global $\max(F)$)};
    \node[stage, below=of red] (tone) {Pass 3: Dynamic color mapping (normalization, colormap, opacity threshold)};
    \node[io, below=of tone] (out) {Final heatmap image + legend + Line Chart};

    \draw[arrow] (raw) -- (merge);
    \draw[arrow] (merge) -- (inst);
    \draw[arrow] (inst) -- (vs);
    \draw[arrow] (vs) -- (acc);
    \draw[arrow] (acc) -- (red);
    \draw[arrow] (red) -- (tone);
    \draw[arrow] (tone) -- (out);

    \draw[arrow] (red.west) -- ++(-5mm,0) |- node[pos=0.25, left, align=right]{async readback} (out.west);
  \end{tikzpicture}
  \caption{End-to-end rendering pipeline. Visible samples are clipped and merged on the CPU, then submitted to a three-pass GPU pipeline. In Pass 1, the vertex shader expands one instanced screen-aligned quad per merged kernel, and the fragment shader performs KDE accumulation via quad rasterization with additive blending into an HDR texture. Pass 2 is a compute shader that performs global-maximum reduction. Pass 3 performs dynamic color mapping. The reduction output is consumed by both the color-mapping shader (in the same frame) and by asynchronous CPU readback for legend scaling.}
  \label{fig:pipeline-overview}
\end{figure}

\begin{figure}[t]
  \centering
  \subfloat[Reference weighted KDE without point merging ($N=100{,}000$ visible samples).]{%
    \includegraphics[width=\linewidth]{figs/full_image.png}%
  }\\[0.5em]
  \subfloat[Screen-space weighted-merge approximation with merge threshold $\tau=\sigma$, reducing the representation to 275 kernels while preserving the dominant density structures.]{%
    \includegraphics[width=\linewidth]{figs/approx_image.png}%
  }
  \caption{Qualitative comparison between full KDE accumulation and the proposed screen-space weighted-merge approximation. The merged result substantially reduces kernel evaluations while maintaining the salient high-density regions and overall distributional morphology. See supplementary material for these visualizations and more examples.}
  \label{fig:kde-full-vs-merged}
\end{figure}

\subsection{Pass 1 --- Accumulation}

The first pass performs KDE directly on the GPU. Each downsampled point is rendered with \emph{instanced drawing}: one draw call emits many instances of a small screen-aligned quad, and each instance splats one Gaussian kernel onto a high-dynamic-range (HDR) floating-point texture using additive blending. Because each fragment simply adds its contribution to the running sum, no explicit synchronization is needed: hardware blending accumulates the KDE density field correctly even when kernels overlap.

The kernel radius is configurable in screen pixels. The quad is sized to cover the full extent of the kernel --- three standard deviations in each direction --- so that effectively all of the kernel mass is captured. Inside the quad, each fragment evaluates the Gaussian as a function of its distance from the point center, scaled by the point's weight. Points outside the current view window are discarded before reaching the rasterizer.

For an isotropic 2D Gaussian, the truncation error outside a disk of radius $R$ is (via radial form of the Gaussian \cite{Bishop2006})
\begin{equation}
\varepsilon_{2}(R,\sigma)=\exp\left(-\frac{R^2}{2\sigma^2}\right).
\end{equation}
Setting $R=3\sigma$ gives $\varepsilon_{2}(3\sigma,\sigma)\approx 0.011109$. That is, about $1.11\%$ omitted mass. For our purposes, this error level is acceptable.

For each merged kernel, we keep its center $\overrightarrow{\boldsymbol{\mu}}\in\mathbb{R}^2$, weight $w$, and covariance $\overrightarrow{\mathbf{P}}\in\mathbb{R}^{2\times 2}$, splatting each as an anisotropic Gaussian footprint in the spirit of covariance-driven Gaussian splatting \cite{10.1145/3592433}. In the vertex shader, we build a Cholesky factor $\overrightarrow{\mathbf{L}}$ (as in Gaussian sampling methods \cite{Bishop2006}) such that
\begin{equation}
\overrightarrow{\mathbf{P}} = \overrightarrow{\mathbf{L}}\overrightarrow{\mathbf{L}}^T,
\end{equation}
with
\begin{equation}
\overrightarrow{\mathbf{P}}=
\begin{bmatrix}
\sigma_x^2 & \sigma_{xy} \\
\sigma_{xy} & \sigma_y^2
\end{bmatrix}
\end{equation}
\begin{equation}
\overrightarrow{\mathbf{L}}=
\begin{bmatrix}
\ell_{00} & 0 \\
\ell_{10} & \ell_{11}
\end{bmatrix}
\end{equation}
\begin{equation}
\ell_{00}=\sqrt{\sigma_x^2}
\end{equation}
\begin{equation}
\ell_{10}=\frac{\sigma_{xy}}{\ell_{00}}
\end{equation}
\begin{equation}
\ell_{11}=\sqrt{\sigma_y^2-\ell_{10}^2}
\end{equation}
Starting from standardized quad coordinates $\overrightarrow{\boldsymbol{z}}\in[-3,3]^2$, the shader computes
\begin{equation}
\overrightarrow{\boldsymbol{x}} = \overrightarrow{\boldsymbol{\mu}} + s\overrightarrow{\mathbf{L}}\overrightarrow{\boldsymbol{z}},
\end{equation}
where $s$ is the user-controlled kernel size (in pixels). This is the explicit linear map that turns an isotropic footprint into an oriented anisotropic ellipse in screen space.

The fragment shader does not need to invert $\overrightarrow{\mathbf{P}}$ per pixel. Instead, it receives $\overrightarrow{\boldsymbol{z}}$ as a varying from the vertex stage, and rasterization hardware linearly interpolates $\overrightarrow{\boldsymbol{z}}$ across the two triangles of the quad. Therefore, each fragment already has the local standardized coordinates needed for the quadratic form
\begin{equation}
d_M^2 = \overrightarrow{\boldsymbol{z}}^T\overrightarrow{\boldsymbol{z}},
\end{equation}
which is equivalent to the Mahalanobis distance in screen coordinates:
\begin{equation}
d_M^2 = (\overrightarrow{\boldsymbol{x}}-\overrightarrow{\boldsymbol{\mu}})^T (s^2\overrightarrow{\mathbf{P}})^{-1}(\overrightarrow{\boldsymbol{x}}-\overrightarrow{\boldsymbol{\mu}}).
\end{equation}
The kernel value is then evaluated as
\begin{equation}
K(\overrightarrow{\boldsymbol{x}})=\exp\left(-\frac{1}{2}d_M^2\right),
\end{equation}
and the weighted contribution $wK(x)$ is additively accumulated into the HDR buffer. This avoids per-fragment matrix inversion and leverages fixed-function interpolation for efficient Gaussian evaluation.

\subsection{Pass 2 --- Reduction}

The second pass finds the maximum accumulated value across the entire HDR texture. A parallel reduction --- dispatched across the texture in tiles --- scans every texel and records the global maximum via an atomic operation.

This maximum is read back to the CPU asynchronously for use in the legend scale. Because the readback is asynchronous, it introduces at most one frame of lag, which is imperceptible in practice. More importantly, the value is available to the dynamic color mapping pass in the same frame via the same GPU-side buffer, ensuring normalization is always up to date.

\subsection{Pass 3 --- Dynamic Color Mapping}

The third pass maps the raw floating-point accumulation values to a perceptual colormap and writes the final output. A fullscreen quad reads from the HDR texture and the maximum value buffer, normalizing each pixel's density to the range $[0, 1]$.

Pixels below a configurable opacity threshold are rendered fully transparent, hiding empty or nearly empty regions and cleanly separating signal from background.

The normalized value is then mapped through a three-stop colormap with a user-adjustable midpoint. The fragment shader interpolates the color gradient from a low-density color through a mid-density color to a high-density color. The midpoint position controls where the perceptual emphasis of the colormap sits: pulling it toward zero compresses the high end and reveals fine structure in sparse regions; pushing it toward one stretches the low end and emphasizes the densest concentrations.

An optional quantization step maps the normalized value to a discrete number of levels before color lookup, producing a contour-map aesthetic that can make density gradients easier to read.


\section{Results and Discussion}
\label{sec:results}

%We presented an interactive density-heatmap pipeline for weighted timeseries data. Three GPU passes --- additive Gaussian accumulation, max-reduction, and dynamic color mapping --- are sufficient for real-time visualization at financial-tick scale. Adaptive weighted-merge downsampling bounds work to display resolution, so performance scales with visible data rather than total dataset size.

\subsection{Benchmark Results}

Table~\ref{tab:merge-benchmark} reports merge-pass benchmark results measured on an Intel Core i5-10400F CPU using Node.js 24.12.0, with resolution $1920\times1080$, $\texttt{sigmaSizePx}=16$, and merge threshold $\tau=1$ pixel.

\begin{table}[t]
  \centering
  \caption{Merge benchmark summary on i5-10400F (Node.js 24.12.0), using $1920\times1080$, $\texttt{sigmaSizePx}=16$, and $\tau=1$.}
  \label{tab:merge-benchmark}
  \begin{tabular}{|r|r|r|r|r|r|}
    \hline
    $N$ & merged points & mean (ms) & p99 (ms) & rme (95\%) & samples \\
    \hline
    1{,}000{,}000 & 1813 & 10.3327 & 12.4025 & $\pm 1.19\%$ & 50 \\
    100{,}000 & 1034 & 1.0075 & 1.3377 & $\pm 0.57\%$ & 497 \\
    10{,}000 & 621 & 0.1217 & 0.1922 & $\pm 0.37\%$ & 4108 \\
    1{,}000 & 332 & 0.0302 & 0.0534 & $\pm 7.70\%$ & 16570 \\
    \hline
  \end{tabular}
\end{table}

\subsection{Complexity}

Per-frame CPU work scales linearly with the number of visible points and runs in a single pass. GPU work for the accumulation pass scales with the number of downsampled points and the kernel radius; the reduction and dynamic color mapping passes scale with screen resolution. In practice, the downsampling limits the GPU instance count to a few thousand even for datasets with hundreds of thousands of points, making the accumulation pass faster than the fixed-cost passes.

\subsection{Mean Squared Error Analysis}

Let $F(x,y)$ denote the reference density field obtained from full KDE (no merging), and let $\hat{F}(x,y)$ denote the field produced by the weighted-merge approximation. Over the heatmap domain $\Omega$ with $M$ pixels, we measure approximation quality with
\begin{equation}
\mathrm{MSE}=\frac{1}{M}\sum_{(x,y)\in\Omega}\left(\hat{F}(x,y)-F(x,y)\right)^2,
\end{equation}
and its scale-normalized form
\begin{equation}
\mathrm{NRMSE}=\frac{\sqrt{\mathrm{MSE}}}{\max_{(x,y)\in\Omega}F(x,y)-\min_{(x,y)\in\Omega}F(x,y)}.
\end{equation}

To estimate these quantities in practice, we use Monte Carlo sampling over the visible domain with the same setup used in the benchmark ($1920\times1080$, $\texttt{sigmaSizePx}=16$, and $\tau=1$). At each iteration, we first evaluate the approximation $\hat{F}$. If $\hat{F}<0.01$, the sample is rejected and a new point is drawn. This acceptance rule avoids over-representing empty background regions that would artificially lower the measured error. For accepted samples, we evaluate both $F$ and $\hat{F}$, accumulate squared error, and compute MSE and NRMSE. We also report relative margin of error (RME, 95\%) for the MSE estimate. % as
%\begin{equation}
%\mathrm{RME}=100\cdot\frac{1.96\,\mathrm{SE}(\mathrm{MSE})}%{\mathrm{MSE}}.
%\end{equation}

Table~\ref{tab:monte-carlo-error} summarizes the measured approximation quality at three dataset sizes.

\begin{table}[t]
  \centering
  \caption{Monte Carlo approximation quality for weighted-merge KDE.}
  \label{tab:monte-carlo-error}
  \begin{tabular}{|r|r|r|r|}
    \hline
    $N$ & merged kernels & NRMSE (\%) & RME (95\%) \\
    \hline
    1{,}000 & 332 & 0.0979\% & $\pm 22.63\%$ \\
    10{,}000 & 621 & 0.0625\% & $\pm 16.40\%$ \\
    100{,}000 & 1034 & 0.0301\% & $\pm 17.48\%$ \\
    \hline
  \end{tabular}
\end{table}

To isolate merge effects, MSE is measured on the linear HDR accumulation buffer before opacity thresholding and colormapping. Increasing threshold $\tau$ reduces submitted kernels and frame time but increases MSE, so $\tau$ acts as a quality-performance knob. Because merging is recomputed in screen space per view, the error is naturally view-dependent: zoomed-out views tolerate more merging, while zoomed-in views converge toward full KDE.

These empirical results support the use of merge thresholds on the order of $\sigma$ to balance quality and performance.

\subsection{Limitations}

The technique requires a GPU that supports floating-point render target formats with additive blending. This is widely supported on desktop hardware but may not be available on all mobile devices.

At very high zoom levels --- sub-millisecond time windows for tick data --- the view contains so few points that the density field degenerates into isolated splats, and the heatmap character of the visualization disappears. At that scale, a plain line or scatter chart is more appropriate.

\section{Conclusions}
\label{sec:conclusions}

We presented an interactive density-heatmap pipeline for weighted timeseries data. Three GPU passes --- additive Gaussian accumulation, max-reduction, and dynamic color mapping --- are sufficient for real-time visualization at financial-tick scale. Adaptive weighted-merge downsampling bounds work to display resolution, so performance scales with visible data rather than total dataset size.

Although we have not conducted formal usability tests, the technique was motivated by the need to provide users with support for analyzing large time series of stock trading and detecting patterns of value and volume over time. After an informal assessment by an expert from a fintech company, the technique is being integrated into their product as an indicator available to their clients. 

\subsection{Future Work}

There are several venues we are envisioning as future work. At first, we will compare our approach with other techniques in terms of performance and error.  

The shape-preserving downsampling strategies (LTTB, RDP) could drive a dedicated line chart pass in parallel with the density heatmap, providing a geometrically accurate price chart without inflating kernel weights. We plan to combine both representations in a single view, making the relationship between price trajectory and activity density (negotiated volume) immediately readable.

The colormap is defined in display color space, while kernel accumulation operates in linear space. Perceptual tone-mapping could be used to improve visual encoding, making information easier to interpret.

Finally, we will assess it across different time series to demonstrate its generalizability. Such time series might contain multiple stocks, different temporal scales, and also different volume distributions.

% conference papers do not normally have an appendix


% trigger a \newpage just before the given reference
% number - used to balance the columns on the last page
% adjust value as needed - may need to be readjusted if
% the document is modified later
%\IEEEtriggeratref{8}
% The "triggered" command can be changed if desired:
%\IEEEtriggercmd{\enlargethispage{-5in}}

% references section

% can use a bibliography generated by BibTeX as a .bbl file
% BibTeX documentation can be easily obtained at:
% http://mirror.ctan.org/biblio/bibtex/contrib/doc/
% The IEEEtran BibTeX style support page is at:
% http://www.michaelshell.org/tex/ieeetran/bibtex/
\bibliographystyle{IEEEtran}
% argument is your BibTeX string definitions and bibliography database(s)
\bibliography{example}
%
% <OR> manually copy in the resultant .bbl file
% set second argument of \begin to the number of references
% (used to reserve space for the reference number labels box)
%\begin{thebibliography}{1}
%
%\bibitem{IEEEhowto:kopka}
%H.~Kopka and P.~W. Daly, \emph{A Guide to \LaTeX}, 3rd~ed.\hskip 1em plus
%  0.5em minus 0.4em\relax Harlow, England: Addison-Wesley, 1999.

%\end{thebibliography}




% that's all folks
\end{document}


